%% ============================================================
%% APPENDICES
%% ============================================================
\appendix

\chapter{Python Model Code}
\label{app:code}

All models are contained in the Python package \texttt{co2whr} and four run
scripts, supplied with this thesis. Running \texttt{python run\_all.py} from
the \texttt{code/} directory regenerates every figure in \texttt{figures/}
and writes every number used in this thesis to
\texttt{results/results.json} and \texttt{results/summary.txt} (about two
minutes on a laptop). Python\,3.9+, CoolProp\,6.x (HEOS backend), NumPy and
Matplotlib are required.

\begin{description}[leftmargin=2em,itemsep=3pt]
\item[\texttt{co2whr/cycles.py}] Steady-state state-point models for the
SRC, sCO$_2$ Brayton and tCO$_2$ Rankine cycles with a recuperator limited
by both streams (Eq.~\ref{eq:recup}); flow exergy and component exergy
balance.
\item[\texttt{co2whr/heatex.py}] Counter-flow heater integration with a
duty-linear gas profile, bisection for the gas outlet temperature at a
10\,K approach, segmented $UA$ (Eq.~\ref{eq:seg_UA}) and two-point LMTD;
recuperator $UA$ with real properties on both sides; cooler $UA$.
\item[\texttt{co2whr/design.py}] Site design points: pinch-limited steam
pressure by golden-section search on $w_{gas}$, recovered power from each
stream, LMTD-versus-segmented comparison.
\item[\texttt{co2whr/costs.py}] Cost correlations of Table~\ref{tab:cost_basis},
CEPCI escalation, currency conversion and total erected cost by component.
\item[\texttt{co2whr/economics.py}] LCOE, NPV, discounted and simple payback,
IRR, Monte Carlo sampling and heat-exchanger breakeven.
\item[\texttt{co2whr/offdesign.py}] Stodola part-load correction with
sliding-pressure control; Marchionni flexibility scoring.
\item[\texttt{co2whr/sites.py}] The three case-study definitions and shared
constants.
\item[\texttt{run\_rq1.py} -- \texttt{run\_rq4.py}, \texttt{run\_all.py}]
One script per research question producing the figures and result tables;
\texttt{run\_all.py} runs all four and writes the results files.
\end{description}

Validation outcomes: SRC model versus NIST Webbook, maximum deviation
0.031\,\% on specific enthalpy; CO$_2$ critical point versus NIST REFPROP,
0.04\,\%; global exergy balance closure below $10^{-12}$\,kJ/kg for all
three cycles; recuperator internal temperature difference positive at every
node for every reported case. The Interim Report model files
(\texttt{tco2\_model.py}, \texttt{sco2\_model.py}, \texttt{src\_model.py},
\texttt{rq2\_sizing.py}, \texttt{rq2\_cost.py}, \texttt{rq3\_economics.py},
\texttt{rq4\_offdesign.py}) are retained in the repository under
\texttt{code/interim\_original/} for comparison; their known errors are
listed in Table~\ref{tab:revisions}.

\chapter{Complete Findings Log}
\label{app:findings}

All thirty findings are listed below with their section reference and key
numerical value.

\begin{longtable}{>{\raggedright\arraybackslash}p{1.0cm}
  >{\raggedright\arraybackslash}p{5.2cm}
  >{\raggedright\arraybackslash}p{8.0cm}}
\toprule
\rowcolor{tableheader}
\textbf{No.} & \textbf{Finding} & \textbf{Key value} \\
\midrule
\endfirsthead
\toprule
\rowcolor{tableheader}
\textbf{No.} & \textbf{Finding} & \textbf{Key value} \\
\midrule
\endhead
\bottomrule
\endlastfoot
\multicolumn{3}{l}{\textit{RQ1 --- Thermodynamic performance (Chapter~\ref{ch:results_rq1})}} \\
F1  & SRC most efficient at every TIT, 150--600\,\textdegree C & SRC 19.1--35.4\,\%; tCO$_2$ 13.3--30.1\,\%; sCO$_2$ 5.3--30.1\,\%; no crossover \\
F2  & sCO$_2$ BWR prohibitive below 300\,\textdegree C & 79.8\,\% at 150\,\textdegree C; 45\,\% at 300\,\textdegree C; 26.9\,\% at 600\,\textdegree C \\
F3  & Recuperator gain bounded by the hot stream & +3.4/+7.0/+10.1 pp at $\varepsilon=0.70$ (interim: +12.3/+16.8/+20.3) \\
F4  & Steam boiler largest single loss; steam still most exergy-efficient per unit work at 400\,\textdegree C & 375\,kJ/kg (56.6\,\%); 0.70 vs 1.09 (tCO$_2$) vs 1.25 (sCO$_2$) kJ/kJ \\
F5  & Recuperator the dominant CO$_2$ irreversibility & 28\,\% (sCO$_2$) and 32\,\% (tCO$_2$) of destruction at 400\,\textdegree C \\
F6  & tCO$_2$ heater pinched at both terminals, bulge mid-exchanger & $\Delta T_{max}$ 41.6/49.0/64.7\,K at 88/145/166\,\textdegree C \\
F7  & Heater inlet sets the floor on recovery; recuperation raises it & Gas out 34--42\,\textdegree C ($\varepsilon=0$) vs 60/101/151\,\textdegree C ($\varepsilon=0.70$) \\
F8  & Pseudo-critical peak sharpens as pressure falls & 4.42\,kJ/kg$\cdot$K at 57\,\textdegree C (12.8\,MPa); 2.62 at 76\,\textdegree C (20\,MPa) \\
F9  & Recuperation adds more to efficiency than to recovered work & $\varepsilon$ 0$\to$0.85: $\eta$ +4.5/+9.3/+13.9 pp; $w_{gas}$ +12/+16/+21\,\% \\
F10 & High $\varepsilon$: recuperator approaches heater duty and the approach limit & 302 vs 399\,kJ/kg at CS3; $\Delta T_{min}$ 9.8\,K at CS2 \\
F11 & Low condensing pressures climatically infeasible & 3.5/4.0/5.0\,MPa condense at 0.2/5.3/14.3\,\textdegree C \\
F12 & Chiller for a 14\,\textdegree C sink would consume most or all of the output & 134/75/63\,\% of net work at COP 4 \\
F13 & Warm sink penalty on sCO$_2$ concentrated at low TIT & 32$\to$40\,\textdegree C inlet: 15.7$\to$9.2\,\% at 200\,\textdegree C \\
\midrule
\multicolumn{3}{l}{\textit{RQ2 --- Sizing and capital cost (Chapter~\ref{ch:results_rq2})}} \\
F14 & Integrated steam pressure far below the sweep value & 0.16/0.91/3.69\,MPa; site $\eta$ 14.4/23.7/30.8\,\% \\
F15 & tCO$_2$ recovers the most power at 200 and 350\,\textdegree C; sCO$_2$ the least everywhere & 723/3418/2186\,kW vs SRC 412/2639/2273 vs sCO$_2$ 340/1988/1280 \\
F16 & CO$_2$ turbines 15--26$\times$ more compact & SP 0.63--0.81 (SRC) vs 0.031--0.042 (CO$_2$) \\
F17 & Steam needs multi-stage turbines above 350\,\textdegree C; CO$_2$ single-stage & VFR 10.5/45/148 (SRC) vs 2.0--3.1 (CO$_2$) \\
F18 & Two-point LMTD overestimates tCO$_2$ heater area & +129/+171/+223\,\% at $\varepsilon=0$ \\
F19 & Steam cheapest per kW at every site & 4{,}444/3{,}071/2{,}598 AUD/kW at 1\,MWe; tCO$_2$ +30--65\,\%; sCO$_2$ 2.3--2.6$\times$ \\
F20 & sCO$_2$ penalty is turbomachinery & 2.8\,MW turbine + 1.8\,MW compressor per MWe at CS2; 37\,\% of equipment \\
F21 & Heat exchangers 35--43\,\% of CO$_2$ equipment; heat rejection 17--44\,\% for all & Cooler largest for 14\,\%-efficient steam at CS2 \\
\midrule
\multicolumn{3}{l}{\textit{RQ3 --- Techno-economics (Chapter~\ref{ch:results_rq3})}} \\
F22 & Steam lowest LCOE at every site & 96/35/31 AUD/MWh; tCO$_2$ 100/53/57; sCO$_2$ 243/94/91 \\
F23 & tCO$_2$ best NPV at 200\,\textdegree C; steam at 350 and 500\,\textdegree C & CS2: 1.71 vs 1.12\,M; CS1: 11.4 vs 10.1\,M; CS3: 12.4 vs 7.5\,M \\
F24 & Discounted payback 3--4\,yr steam, 7\,yr tCO$_2$ at hot sites; 10--11\,yr at CS2 & IRR 28/36\,\% (SRC), 17/18\,\% (tCO$_2$), 12\,\% at CS2 \\
F25 & Economics follow efficiency $\times$ recovery depth & sCO$_2$ at CS1: $-1.4$ pp but $-25$\,\% power; tCO$_2$ at CS2: +1.2 pp but +75\,\% power \\
F26 & Steam robust at 350/500\,\textdegree C; CS2 a one-in-three risk & $P(\mathrm{NPV}>0)$ 99.8/100\,\% (SRC), 90/91\,\% (tCO$_2$), 72/68\,\% at CS2 \\
F27 & Cheaper HX rescues tCO$_2$ only at 200\,\textdegree C & 11\,\% cut at CS2; 95\,\% at CS1; unreachable at CS3 and for sCO$_2$ \\
\midrule
\multicolumn{3}{l}{\textit{RQ4 --- Off-design (Chapter~\ref{ch:results_rq4})}} \\
F28 & tCO$_2$ retains most efficiency at deep part load; sCO$_2$ least & 56--67\,\% vs 38--66\,\% (SRC) vs 41--47\,\% (sCO$_2$) at 20\,\% load \\
F29 & Steam part-load behaviour depends on design pressure & 66\,\% retention at CS3 (3.7\,MPa) vs 38\,\% at CS2 (0.16\,MPa) \\
F30 & sCO$_2$ 80\,\%-load gain at CS2 is a pressure-ratio effect & Retention 1.006; disappears at 350/500\,\textdegree C \\
\end{longtable}

\chapter{AI Assistance Log}
\label{app:ai}

%% STUDENT: this log must match your own records. Add, remove or reword
%% entries so that it is complete and accurate before submission.
This log records the use of generative AI tools in this thesis, as required
by UQ policy~\cite{UQ2026AI}. Each entry states the section affected, the
tool, what was requested, what was received and how it was used.

\begin{longtable}{>{\raggedright\arraybackslash}p{2.2cm}
  >{\raggedright\arraybackslash}p{1.6cm}
  >{\raggedright\arraybackslash}p{4.3cm}
  >{\raggedright\arraybackslash}p{6.0cm}}
\toprule
\rowcolor{tableheader}
\textbf{Section} & \textbf{Tool} & \textbf{Request} & \textbf{Output and use} \\
\midrule
\endfirsthead
\toprule
\rowcolor{tableheader}
\textbf{Section} & \textbf{Tool} & \textbf{Request} & \textbf{Output and use} \\
\midrule
\endhead
\bottomrule
\endlastfoot
Ch.~1--2 & Claude & Draft background and literature-review paragraphs from my reading notes and reference list & Draft prose; checked against the primary sources, reorganised and rewritten where needed \\
Ch.~1--3 & Claude & \LaTeX{} formatting of definitions, equations, tables and the methodology flow diagram & Typesetting only; content from my proposal and interim report \\
Model code (interim) & Copilot & Autocomplete while writing \texttt{tco2\_model.py} and the other interim scripts & CoolProp/NumPy/Matplotlib syntax completions; all model logic written by me \\
Model code (final) & Claude & Review my seven interim Python files and their printed outputs & Identified: recuperator limited by the cold stream only (second-law violation); gas profile linear in node index in the pinch diagrams; cancelling temperature-difference expression for the sCO$_2$ heater and ad hoc boiler $\Delta T$ in \texttt{rq2\_cost.py}; condenser cost defined but unused; turbine coefficients not traceable to the cited paper; simple payback; Monte Carlo not varying source temperature as claimed \\
Model code (final) & Claude & Restructure the scripts into one package implementing the corrections, a consistent heater integration, the cost basis of Table~\ref{tab:cost_basis}, fixed-source economics, discounted payback and breakeven & \texttt{co2whr} package and run scripts (Appendix~\ref{app:code}). I ran the package, checked balance closures, verified every correlation against Weiland et al.~(2019), Turton et al.~(2012) and DOE (2012), and compared outputs with my interim results \\
Sec.~\ref{sec:method_cost} & Claude & Locate and verify the Weiland et al.~(2019) correlation coefficients & Coefficients and validity ranges from the NETL/DOE source; used in Table~\ref{tab:cost_basis} after my own check of the paper \\
Figures & Claude & Matplotlib scripts for all figures from the model outputs & Figures regenerated by \texttt{run\_all.py}; captions written with the chapter text \\
Ch.~4--9 & Claude & Draft the results, discussion and conclusion chapters from \texttt{results/summary.txt} and my notes on interpretation & Draft text with finding boxes; I checked every number against the results file, revised the interpretation and approved the final wording \\
Abstract, Exec.\ Summary, Declaration & Claude & Draft from the final chapters & Revised and approved by me \\
Conference slides & Claude & Build and revise the presentation from the thesis content and supervisor feedback & Slides and speaker notes; presented by me on 8~October~2026 \\
\end{longtable}
